\documentclass[11pt]{article}
\usepackage{amsmath,amssymb,amsthm,geometry}
\geometry{margin=25mm}\newtheorem{theorem}{Theorem}
\newcommand{\Tr}{\operatorname{Tr}}
\begin{document}
\title{Twenty-six qubits: certified exact sixth-order anticoherent state}
\author{Rational polynomial equations and an exact contraction certificate}
\date{5 November 2026}\maketitle
\begin{theorem}
An exactly defined algebraic pure symmetric twenty-six-qubit state is sixth-order
anticoherent. Together with the independent seventh-order exclusion in
\texttt{../n26-t7/proof.tex}, this proves
\[\boxed{t_{\max}(26)=6.}\]
\end{theorem}
\begin{proof}
This is an exact implicit algebraic construction. Its Dicke coefficients are
defined by a unique root of rational quadratic equations in a specified rational
box. The displayed decimal components are approximations to that state; no
closed-form radical expressions for its components are claimed.

\paragraph{Rational reduced-state equations.}
Write $n=26$, $t=6$, and let the 27 positive integer scales $d_r$ be the powers
of two recorded in \texttt{certificate.json}. For 54 real variables $x_0,\ldots,x_{53}$,
put
\[
z_r=x_r+ix_{27+r},\qquad
c_r=\frac{\sqrt{\binom{26}{r}}}{d_r}z_r,\qquad
|\psi(x)\rangle=\sum_{r=0}^{26}c_r|D_{26}^{(r)}\rangle.
\]
The Dicke states are normalized. The Schmidt matrix for the six-qubit reduction is
\[
M_{a\ell}=\frac{\sqrt{\binom6a\binom{20}\ell}}{d_{a+\ell}}z_{a+\ell},
\qquad0\le a\le6,\quad0\le\ell\le20.
\]
Consequently $\rho_6=MM^\dagger=I_7/7$ is equivalent to 49 real rational
quadratic equations $F(x)=0$. The seven diagonal equations are
\[
F_{aa}=\sum_{\ell=0}^{20}\frac{\binom{20}\ell}{d_{a+\ell}^2}
\bigl(x_{a+\ell}^2+x_{27+a+\ell}^2\bigr)
-\frac1{7\binom6a}=0.
\]
For $0\le a<b\le6$, the other 42 equations are the real and imaginary parts of
\[
F_{ab}=\sum_{\ell=0}^{20}\frac{\binom{20}\ell}{d_{a+\ell}d_{b+\ell}}
z_{a+\ell}\overline{z_{b+\ell}}=0.
\]
The equation ordering is seven diagonal entries, followed by real and imaginary
parts of the upper entries in lexicographic order. Taking the trace of
$\rho_6=I_7/7$ proves $\sum_r|c_r|^2=1$, so no separate normalization premise is needed.

\paragraph{An exact isolating box.}
Let $x^0\in\mathbb Q^{54}$ be the center in the certificate. Fix coordinates
$\{11,13,28,42,48\}$ to their rational central values, and let $y\in\mathbb R^{49}$
contain the other coordinates in the recorded order. This restriction is a
construction choice, not a condition imposed on arbitrary AC states.
The certificate supplies a rational $49\times49$ matrix $A$, whose entries
are integers divided by $2^{256}$, and radius
\[r=10^{-65}.
\]
For the square Jacobian $J_0=DF(y^0)$, exact rational arithmetic verifies
\[
\eta=\|AF(y^0)\|_\infty<10^{-67},\qquad
\beta=\|I-AJ_0\|_\infty<10^{-76},\qquad
\|A\|_\infty<17000.
\]
All equations are quadratic. Each contains two quadratic monomials per
$\ell$, and each monomial contributes at most twice its absolute coefficient
to the sum of absolute Hessian entries. Since $d_r\ge1$, the universal bound
\[
H\le4\sum_{\ell=0}^{20}\binom{20}\ell=2^{22}
\]
implies $\|DF(y)-J_0\|_\infty\le H\|y-y^0\|_\infty$.
Thus on the closed box $B=\{y:\|y-y^0\|_\infty\le r\}$, the map
\[T(y)=y-AF(y)
\]
has Lipschitz constant at most
\[
q=\beta+\|A\|_\infty H r<10^{-54}<1.
\]
Moreover $\|T(y)-y^0\|_\infty\le\eta+qr<r$, as the verifier checks using
the full rational bounds, not their rounded displays. Banach's theorem yields
a unique fixed point $y^*\in B$.

Because $\beta<1$, the matrix $AJ_0$ is invertible, and hence $A$ is invertible.
Therefore $AF(y^*)=0$ implies $F(y^*)=0$. The same derivative bound makes
$DF(y^*)$ invertible. The root is thus an isolated nonsingular zero of a square
rational polynomial system, so its coordinates are real algebraic numbers.
Together with the five fixed rational coordinates, it defines a precise exact
state $|\psi(x^*)\rangle$ satisfying $\rho_6=I_7/7$.

\paragraph{Meaning of the decimal components.}
The scales satisfy $d_r^2\le\binom{26}{r}<4d_r^2$.
Each real or imaginary Dicke component of the exact state is therefore within
$2\cdot10^{-65}$ of the exact center conversion
$\sqrt{\binom{26}{r}}(x_r^0+ix_{27+r}^0)/d_r$.
Binary floating-point components and plotted Majorana coordinates are numerical
renderings. The exact state is specified by the rational equations, fixed
coordinates and isolating box, rather than by those floating values.

\paragraph{Maximum order.}
The separate exact three-copy partial-transpose certificate excludes every
seventh-order pure symmetric twenty-six-qubit state. Nesting excludes all
higher orders. The certified sixth-order state therefore has exactly order six
and attains the global maximum.
\end{proof}
\paragraph{Independent reproduction.}
Run \texttt{python3 verify\_independently.py}; only Python's standard library is
needed. It reconstructs the reduced-state equations, obtains the Jacobian by
exact central differences of the quadratic equations, and verifies every
rational contraction inequality. No discovery library, numerical convergence
claim or floating-point tolerance is a premise. \texttt{build.py} is optional
discovery code and accepts \texttt{--source}; it is not imported by the verifier.
\end{document}
